\section*{4. イオンの名称とイオン式}
\begin{enumerate}
    \item 原子が電子を失って＋の電気を帯びた粒子を（ \textcolor{red}{\textbf{陽}} ）イオン、電子を受け取って－の電気を帯びた粒子を（ \textcolor{red}{\textbf{陰}} ）イオンという。
    \item 次のイオンのイオン式を書きなさい。
    
    \vspace{0.2cm}
    \hspace*{-1.5cm}
    {\renewcommand{\arraystretch}{1.5}
    \begin{tabular}{llllll}
    水素イオン&（ \textcolor{red}{$\mathbf{H^+}$} ） & 酸化物イオン&（ \textcolor{red}{$\mathbf{O^{2-}}$} ） & 硫酸イオン&（ \textcolor{red}{$\mathbf{SO_4^{2-}}$} ） \\
    カリウムイオン&（ \textcolor{red}{$\mathbf{K^+}$} ） & カルシウムイオン&（ \textcolor{red}{$\mathbf{Ca^{2+}}$} ） & バリウムイオン&（ \textcolor{red}{$\mathbf{Ba^{2+}}$} ） \\
    塩化物イオン&（ \textcolor{red}{$\mathbf{Cl^-}$} ） & 水酸化物イオン&（ \textcolor{red}{$\mathbf{OH^-}$} ） & 硫化物イオン&（ \textcolor{red}{$\mathbf{S^{2-}}$} ） \\
    鉄イオン&（ \textcolor{red}{$\mathbf{Fe^{2+}}$} ） & 銅イオン&（ \textcolor{red}{$\mathbf{Cu^{2+}}$} ） & 亜鉛イオン&（ \textcolor{red}{$\mathbf{Zn^{2+}}$} ） \\
    マグネシウムイオン&（ \textcolor{red}{$\mathbf{Mg^{2+}}$} ） &&& & \\
    \end{tabular}}
    \vspace{0.2cm}

    \item 複数の原子が結びついてできているイオンを（ \textcolor{red}{\textbf{多原子イオン}} ）という。
    \item 次のア～カから、(3)に該当するものを \textbf{すべて} 選びなさい。\\[0.2cm]
    ア. アンモニウムイオン \hspace{0.5cm} イ. カリウムイオン \hspace{0.5cm} ウ. 水酸化物イオン \hspace{0.5cm} エ. 塩化物イオン \hspace{0.5cm} オ. 硝酸イオン \hspace{0.5cm} カ. 硫酸イオン\\[0.2cm]
    答（ \textcolor{red}{\textbf{ア、ウ、オ、カ}} ）
\end{enumerate}
